\documentclass[11pt,a4paper]{article}
\usepackage[T1]{fontenc}
\usepackage[margin=27mm]{geometry}
\usepackage{amsmath,amssymb,amsthm,mathtools}
\usepackage{booktabs,array}
\usepackage[expansion=false]{microtype}
\usepackage[hidelinks]{hyperref}
\usepackage{url}
\makeatletter
\g@addto@macro\UrlBreaks{\do\0\do\1\do\2\do\3\do\4\do\5\do\6\do\7\do\8\do\9}
\makeatother
\setlength{\emergencystretch}{3em}
\allowdisplaybreaks[2]
\renewcommand{\labelitemi}{$\bullet$}
\numberwithin{equation}{section}
\newtheorem{maintheorem}{Theorem}
\renewcommand{\themaintheorem}{\Alph{maintheorem}}
\newtheorem{maincorollary}[maintheorem]{Corollary}
\newtheorem{theorem}{Theorem}[section]
\newtheorem{proposition}[theorem]{Proposition}
\newtheorem{lemma}[theorem]{Lemma}
\newtheorem{corollary}[theorem]{Corollary}
\theoremstyle{definition}
\newtheorem{definition}[theorem]{Definition}
\theoremstyle{remark}
\newtheorem{remark}[theorem]{Remark}
\newcommand{\Q}{\mathbb Q}
\newcommand{\R}{\mathbb R}
\newcommand{\Pone}{\mathbb P^1}
\newcommand{\trdeg}{\operatorname{trdeg}}
\newcommand{\crat}{\operatorname{cr}}
\newcommand{\Mat}{\operatorname{Mat}}
\newcommand{\Span}{\operatorname{span}}
\newcommand{\diag}{\operatorname{diag}}
\newcommand{\Der}{\operatorname{Der}}
\newcommand{\eps}{\varepsilon}
\newcommand{\cA}{\mathcal A}
\newcommand{\cB}{\mathcal B}
\newcommand{\cG}{\mathcal G}
\newcommand{\cH}{\mathcal H}
\newcommand{\cR}{\mathcal R}
\newcommand{\HH}{H_0}
\newcommand{\ancfile}[1]{\texttt{\detokenize{#1}}}
\hypersetup{pdftitle={Concyclic Euclidean Ramsey sets with three-dimensional moduli},pdfauthor={Jizhou Guo}}
\title{Concyclic Euclidean Ramsey sets\protect\\ with three-dimensional moduli}
\author{Jizhou Guo\\\texttt{mitsuha2021b@gmail.com}}
\date{October 23, 2026}
\begin{document}
\maketitle
\begin{abstract}
Let $A$ be a finite set of at least seven points on a circle, and let $M$ be
the field generated over $\Q$ by its four-point cross ratios. When
$\trdeg_\Q M=3$, the set is Euclidean Ramsey precisely when it is a union of
two groups, each containing at least three points and having only algebraic
internal cross ratios. For seven points this is equivalent to the existence
of an algebraic four-point cross ratio. We prove the necessary implication
by a second-order tensor obstruction and a classification of simultaneous
totally isotropic three-dimensional subspaces. The sufficient implication
is an application of OpenAI's Theorem~1.1 and is also covered by Strinz's
two-group theorem. The criterion decides the three-parameter example left
open in Appendix~E of P\'alv\"olgyi's paper: it is non-Ramsey. This negative
conclusion, and the necessary implication above, use neither direction of
OpenAI's classification; the needed tensor necessity is proved here.
An explicit elliptic family, whose Ramsey conclusion uses the sufficient
direction of that classification, has one-dimensional moduli and no
algebraic four-point cross ratio. It illustrates the role of the dimension
hypothesis; other examples of this phenomenon are already public.
\end{abstract}

\section{Statements and background}\label{sec:introduction}

A finite nonempty set $A$ in a Euclidean space is \emph{Ramsey} if, for
every integer $r\ge2$, there exists $N\ge1$ such that every colouring
$\R^N\to\{1,\ldots,r\}$ contains a monochromatic congruent copy of $A$.
The scale is fixed, and the colouring has no regularity assumption.
Ramsey-ness depends only on the distance table, is preserved by nonzero
similarities, and passes to subsets. The assertions concerning circle
sets below apply to circles of positive radius in any Euclidean space:
we first choose an isometric planar representative and then use a
similarity to the unit circle.

Use homogeneous parameters $z=[s:t]\in\Pone(\R)$ for the unit circle:
\begin{equation}\label{eq:circle-map}
 a([s:t])=\left(\frac{s^2-t^2}{s^2+t^2},
                       \frac{2st}{s^2+t^2}\right).
\end{equation}
For nonzero representatives put $[z_i,z_j]=s_it_j-t_is_j$. For four
distinct indices define
\begin{equation}\label{eq:cross-ratio}
 \crat(i,j;k,l)=
 \frac{[z_i,z_k][z_j,z_l]}{[z_i,z_l][z_j,z_k]}.
\end{equation}
All denominators are nonzero. This quantity is unchanged by rescaling
representatives or applying a common element of $\mathrm{PGL}_2(\R)$.
Changing the projective parametrization of the circle therefore does
not change the following field.

\begin{definition}\label{def:moduli}
The \emph{projective moduli field} of a labelled circle set $A$ is
\[
 M(A)=\Q\bigl(\crat(i,j;k,l):i,j,k,l\text{ distinct}\bigr).
\]
A group of at least three of its points is \emph{projectively algebraic}
if all its internal four-point cross ratios are algebraic over $\Q$.
For a group of three points this condition is vacuous.
\end{definition}

Equivalently, a projectively algebraic group is a projective image of a
finite subset of $\Pone(\overline{\Q}\cap\R)$. Indeed, sending three of
its points to $0,1,\infty$ makes every remaining parameter a cross ratio
or a fractional linear transform of one. The same observation shows
that $M(A)$ is generated by the remaining $n-3$ parameters after such a
normalization. In particular, the moduli degree refers to the normalized
parameters, rather than the absolute coordinate field.

\begin{maintheorem}\label{thm:main}
Let $A$ consist of $n\ge7$ distinct points on a circle and suppose that
$\trdeg_\Q M(A)=3$. Then the following are equivalent:
\begin{enumerate}
\item $A$ is Ramsey;
\item its indices have a partition $I\sqcup J$, with $|I|,|J|\ge3$,
such that each group is projectively algebraic.
\end{enumerate}
The implication \textup{(2)}$\Rightarrow$\textup{(1)} uses the sufficient
direction of OpenAI's Theorem~1.1~\cite{openai2026}; it is also a circle
case of Strinz's two-group theorem~\cite[Theorem C2]{strinz2026spherical}.
The implication \textup{(1)}$\Rightarrow$\textup{(2)} is proved here
without either direction of OpenAI's theorem.
\end{maintheorem}

\begin{maincorollary}\label{cor:seven}
For seven distinct circle points with $\trdeg_\Q M(A)=3$, the set is
Ramsey if and only if some four-point cross ratio is algebraic over $\Q$.
The sufficient implication uses OpenAI's Theorem~1.1, sufficient
direction; the necessary implication uses neither direction.
\end{maincorollary}

In a partition of seven indices into groups of size at least three the
sizes are $3$ and $4$. For a four-point group all orderings of its cross
ratio are fractional linear transforms of one of them. Thus the
equivalence in Corollary~\ref{cor:seven} follows immediately from
Theorem~\ref{thm:main}. Its sufficient implication holds without a
moduli-degree hypothesis, as already follows from Strinz's results.

\begin{maincorollary}\label{cor:palvolgyi}
Let $x,y,z$ be algebraically independent real numbers over $\Q$. The
seven circle points with finite parameters
\begin{equation}\label{eq:appendix-family}
 -2,\quad 0,\quad 2,\quad x,\quad y,\quad -4-x-2y,\quad z
\end{equation}
are non-Ramsey. This conclusion uses neither direction of OpenAI's
Theorem~1.1.
\end{maincorollary}

This is the family discussed after Theorem~A.8 in Appendix~E of
P\'alv\"olgyi~\cite{palvolgyi2026}. His parameters are
$s_1=x$, $s_2=y$, $s_3=-4-s_1-2s_2$, and $s_4=z$. That discussion
leaves the Ramsey status open and imposes a neighbourhood condition for
the failure of its derivation-certificate method. The corollary above
has no neighbourhood condition. Section~\ref{sec:appendix-example}
gives a short cross-ratio proof and an independent finite second-order
certificate.

\begin{maintheorem}\label{thm:elliptic}
For each transcendental real $X>\sqrt[3]{2}$ and either choice
$Y=\pm\sqrt{X^3-2}$, the recurrence in
\eqref{eq:elliptic-map} produces seven distinct circle parameters
$X_0,\ldots,X_6$. Their circle set is Ramsey, its moduli field has
transcendence degree one, and all its four-point cross ratios are
transcendental. The Ramsey assertion uses the sufficient direction of
OpenAI's Theorem~1.1.
\end{maintheorem}

This example shows that the degree-three hypothesis in the necessary
implication of Corollary~\ref{cor:seven} cannot simply be omitted.
The existence of Ramsey seven-point sets with no algebraic cross ratio
is already exhibited by Strinz~\cite[Section~7, ``Non-two-group Ramsey
sets'']{strinz2026spherical}, using different and simpler rational
families, including
\[
 \{0,\tfrac12,\tfrac13,t,t+1,2t,1/t\}
 \qquad(t\text{ real and transcendental}).
\]
That positive conclusion also uses the sufficient direction of
OpenAI's Theorem~1.1.
The elliptic example is included as a concrete boundary construction.

\subsection*{Public results used for context}
OpenAI's \emph{A classification of finite Euclidean Ramsey
configurations}, dated September~23, 2026, states a necessary and
sufficient tensor condition in Theorem~1.1~\cite{openai2026}.
We use the release at commit
\texttt{adc7f1241b42e322a6451854ab7e4b4c146bf78a}, result~172.
The corresponding Lean development includes
\ancfile{fieldCriterion_ramsey} and \ancfile{classification}
in \ancfile{OAI/Combinatorics/EuclideanRamsey/Main.lean}
at that commit~\cite{openai2026lean}. Its source supplies both
directions; no Lean kernel replay is part of the present note.

Jones's October~7, 2026 note gives an explicit tensor certificate for
any six concyclic points~\cite[Section~2]{jones2026}. Composed with
the sufficient direction of OpenAI's Theorem~1.1, this proves that every circle set with
at most six points is Ramsey. Strinz's October~15, 2026 release states
this conclusion as Theorem~B(a) and supplies a Lean development
composing the six-point construction with the sufficient direction of
OpenAI's Theorem~1.1
\cite{strinz2026spherical}. Thus the six-point result is background
here.

P\'alv\"olgyi's September~20, 2026 preprint states in Theorem~A.3 that,
for any fixed distinct real $t_1,t_2,t_3$, if $t_4,\ldots,t_7$ are
algebraically independent over $\Q(t_1,t_2,t_3)$, the seven circle
points are non-Ramsey~\cite{palvolgyi2026}. That negative result has a
derivation and colouring proof and uses neither direction of OpenAI's
Theorem~1.1. Its applicability is stronger than independence of
all seven absolute parameters, but does not cover
\eqref{eq:appendix-family}.

Strinz's M\"obius-invariance theorem, released October~13, 2026,
states that a finite spherical set is Ramsey if and only if its image
under a M\"obius transformation of its containing sphere is Ramsey
\cite[Theorem C]{strinz2026mobius}. Its published argument uses both
directions of OpenAI's Theorem~1.1.
The October~15 spherical note gives two further sufficient conditions:
two M\"obius images of finite algebraic subsets of a sphere have Ramsey
union (Theorem~C2), and an arbitrary finite algebraic circle set together
with at most three arbitrary circle points is Ramsey (Theorem~B(b))
\cite{strinz2026spherical}. Both positive conclusions use the sufficient
direction of OpenAI's Theorem~1.1. Their circle case explains the sufficient halves
of Theorem~\ref{thm:main} and Corollary~\ref{cor:seven}. The explicit
certificate is recalled in Section~\ref{sec:positive}.

For comparison, on $S^{d-1}$ with $d\ge3$, algebraically independent
stereographic parameters give the generic threshold
$\binom{d+2}{2}$: fewer points are Ramsey and that many or more are
non-Ramsey. Strinz's ledger SR-C3 and SR-C29 records the positive half
as a consequence of OpenAI's Proposition~7.3, using the sufficient
direction of its Theorem~1.1; the negative boundary is P\'alv\"olgyi's Theorem~A.12 and
uses neither direction~\cite{strinz2026spherical,palvolgyi2026}.
On the circle, the relevant generic negative size is seven, and the
classification below concerns the degree-three locus within its moduli.

The proofs occupy Sections~\ref{sec:tensors}--\ref{sec:classification}.
Section~\ref{sec:appendix-example} treats Corollary~\ref{cor:palvolgyi},
and Section~\ref{sec:elliptic} proves Theorem~\ref{thm:elliptic}.
The ancillary files record exact finite checks; the structural
necessity argument itself is given over the full moduli field.

\section{The tensor criterion and its necessary direction}\label{sec:tensors}

For a field $F\subset\R$ write
\[
 \cB_F=F\otimes_\Q F,\qquad L(x)=x\otimes1,\qquad R(x)=1\otimes x,
 \qquad m(x\otimes y)=xy.
\]
Every nonzero pure factor $L(x)$ or $R(x)$ is a unit. The ring
$\cB_F$ need not be a domain. If $a_1,\ldots,a_n\in F^d$, put
$p_i=(1,a_i)^T$. A \emph{tensor certificate} is a matrix
$P\in\Mat_{d+1}(\cB_F)$ with
\begin{equation}\label{eq:tensor-criterion}
 L(p_i)^TPR(p_i)=0\quad(1\le i\le n),\qquad
 m(P_{\alpha\beta})=\delta_{\alpha\beta}
                   \quad(1\le\alpha,\beta\le d).
\end{equation}
The constant row and column have no prescribed multiplied values.
OpenAI's Theorem~1.1 says that a distinct configuration of at least two
points spanning $\R^d$ is Ramsey if and only if
\eqref{eq:tensor-criterion} holds in its coordinate field
\cite{openai2026}. The sufficient direction is the input for all
positive conclusions of this note. To specify the dependence of the
negative conclusions, we give the elementary necessity argument
separately; it is the tensor separation and colouring mechanism of
\cite[Section~2]{openai2026}, with the required descent made explicit.

\begin{lemma}[Descent]\label{lem:descent}
Let $F\subset E\subset\R$, and suppose the coordinates lie in $F$.
A tensor certificate over $E\otimes_\Q E$ descends to one over
$F\otimes_\Q F$. The same assertion holds if the whole multiplied
matrix is prescribed to be a fixed matrix over $F$.
\end{lemma}
\begin{proof}
Choose an $F$-basis $(e_\lambda)$ of $E$. There is a free-module
decomposition
\[
 E\otimes_\Q E=
 \bigoplus_{\lambda,\mu}(F\otimes_\Q F)(e_\lambda\otimes e_\mu).
\]
Expand the given matrix as a finite sum
$P_E=\sum(e_\lambda\otimes e_\mu)P_{\lambda\mu}$.
Since its evaluation coefficients lie in $F\otimes_\Q F$, each
$P_{\lambda\mu}$ has zero evaluations. If $H_{\lambda\mu}=m(P_{\lambda\mu})$,
the required entries of the target matrix $H$ satisfy
\[
 H=\sum_{\lambda,\mu}e_\lambda e_\mu H_{\lambda\mu},
\]
where only the prescribed entries are included if appropriate.
This is a finite linear system over $F$, soluble over $E$.
Row reduction over $F$ shows that it has a solution
$c_{\lambda\mu}\in F$. Then
$\sum L(c_{\lambda\mu})P_{\lambda\mu}$ is the required certificate.
\end{proof}

\begin{proposition}[Independent necessity]\label{prop:necessity}
If a distinct, affinely spanning configuration $a_1,\ldots,a_n\in\R^d$
has no tensor certificate in its coordinate field, it is non-Ramsey.
In every ambient dimension there is an avoiding colouring with at most
$(2n+1)^n$ colours. This statement uses neither direction of OpenAI's
Theorem~1.1.
\end{proposition}
\begin{proof}
Descent permits us to work with $\R\otimes_\Q\R$. Regard
$W_*=\R\times\R^d$ as a $\Q$-vector space, and write
$e_i(v,u)=v+a_i\cdot u$. On the subspace $T(W_*)$ of symmetric
tensors in $W_*\otimes_\Q W_*$ define
\[
 \Psi(T)=\bigl((e_i\otimes e_i)(T)\bigr)_{i=1}^n\oplus G(T),
 \qquad
 G\bigl((v,u)\otimes(v',u')\bigr)=u(u')^T.
\]
Here $G(T)$ is a real symmetric matrix. Distributing the finite
coordinate sums identifies these tensors with matrices over
$\R\otimes_\Q\R$. Under this identification, the first components
are the tensor evaluations and the last component is the multiplied
spatial block. A certificate can always be made symmetric for the
combined transpose and tensor flip by averaging. Hence
$y_*=(0,\ldots,0,I_d)$ is outside the image of $\Psi$.

Separation in the quotient $\Q$-vector space gives a $\Q$-linear
functional vanishing on $\operatorname{im}\Psi$ and taking $y_*$ to
one: choose a basis containing the nonzero class of $y_*$.
Write its components as
$\ell_i:\R\otimes_\Q\R\to\R$ and
$\Lambda:\operatorname{Sym}_d(\R)\to\R$, and put
$h_i(x)=\ell_i(x\otimes x)$. On a pure square this gives
\begin{equation}\label{eq:separation-functions}
 \sum_i h_i(v+a_i\cdot u)+\Lambda(uu^T)=0,
 \qquad \Lambda(I_d)=1.
\end{equation}
Taking $u=0$ shows $\sum_i h_i(x)=0$ for all real $x$.

In $\R^N$ define $H_i(z)=\sum_{j=1}^N h_i(z_j)$.
Every congruent copy has the form $b_i=t+Aa_i$ with $A^TA=I_d$:
the distance identities identify the Gram matrices of the spanning
difference vectors, giving a linear isometric embedding and then a
translation. If $N<d$ there is no such copy. Apply
\eqref{eq:separation-functions} to each coordinate row of $A$ and
add. Additivity of $\Lambda$ gives
\[
 \sum_iH_i(z)=0\quad(z\in\R^N),\qquad
 \sum_iH_i(b_i)=-\Lambda(A^TA)=-1.
\]
No real scalar is pulled through the merely $\Q$-linear functional.

Partition the residues modulo $2$ of each $H_i(z)$ into $2n+1$
equal half-open intervals of $[0,2)$. Their $n$ labels colour $z$.
In a monochromatic copy,
$H_i(b_i)-H_i(b_1)=2k_i+e_i$ with $k_i\in\mathbb Z$ and
$|e_i|<2/(2n+1)$. Summing gives
$-1=2\sum_i k_i+\sum_i e_i$, where
$|\sum_i e_i|<2n/(2n+1)<1$. An odd integer cannot have absolute
value less than one. This contradiction proves the assertion.
\end{proof}

\subsection{Full normalization on the unit circle}
Put
\[
 h(s,t)=(s^2+t^2,s^2-t^2,2st)^T,\qquad
 \HH=\diag(-1,1,1).
\]
The components of $h$ form a basis of the binary quadratic forms over
any $\Q$-algebra, and
\begin{equation}\label{eq:circle-kernel}
 h(s,t)^T\HH h(u,v)=-2(sv-tu)^2.
\end{equation}
Every bihomogeneous polynomial of bidegree $(2,2)$ over $\cB_F$
therefore has a unique representation $h(s,t)^TPh(u,v)$.
For finite parameters, $h(1,t)^TPh(1,u)$ similarly represents every
polynomial of degree at most two in each variable. These are changes
of basis in free modules and remain valid when $\cB_F$ has zero divisors.

If at least three distinct unit circle points have a certificate,
replace it by $(P+P^*)/2$, where
$P^*_{\alpha\beta}=\tau(P_{\beta\alpha})$ and
$\tau(x\otimes y)=y\otimes x$. Its multiplied matrix is symmetric.
Ordinary evaluation of \eqref{eq:tensor-criterion} gives
\[
 m(P_{00})+1+2m(P_{01})a_{i1}+2m(P_{02})a_{i2}=0.
\]
Three distinct circle points are not collinear. This affine function
must vanish identically, whence
\begin{equation}\label{eq:full-normalization}
 m(P)=\HH.
\end{equation}
Thus prescribing the full matrix here loses no possible certificate.

\subsection{Projective transport without a Ramsey input}\label{sec:transport}
For $A\in\mathrm{GL}_2(\R)$, let $C_A$ be its action on binary
quadratic forms, so $h(Az)=C_Ah(z)$. Equation
\eqref{eq:circle-kernel} and $[Az,Aw]=\det(A)[z,w]$ give
\[
 C_A^T\HH C_A=\kappa\HH,\qquad \kappa=(\det A)^2.
\]
In a common real field containing the coordinates and the entries of
$A$, a certificate with \eqref{eq:full-normalization} transports to
\begin{equation}\label{eq:transport}
 P'=L(\kappa)L(C_A^{-T})PR(C_A^{-1}).
\end{equation}
Indeed the new augmented circle vectors are $C_Ap_i$ divided by their
constant coordinates. These coordinates equal
$\|Az_i\|^2/\|z_i\|^2>0$ and are units in the pure factors.
The evaluations therefore vanish and $m(P')=\HH$.
Lemma~\ref{lem:descent} descends $P'$ to the image coordinate field.
The inverse transformation gives the converse. Thus the existence of
a tensor certificate for a circle set of at least three points is
projectively invariant by algebra alone. When this fact is used to
prove a positive Ramsey conclusion, the sufficient direction of
OpenAI's Theorem~1.1 is still required. Its use in a negative proof
requires only Proposition~\ref{prop:necessity}.

\section{A second-order condition over the moduli field}\label{sec:second-order}

Normalize three distinct parameters to $0,1,\infty$. Choose a rational
$\alpha$ different from every finite normalized parameter, including
$0$ and $1$, and apply
$r\mapsto1/(r-\alpha)$. Every parameter is now finite and the first
three are distinct rational numbers. This operation and its inverse
are rational over $\Q$, so the field of all finite parameters remains
the moduli field $F=M(A)$. The unit circle coordinates have this same
field, since
\[
 a(t)=\left(\frac{1-t^2}{1+t^2},\frac{2t}{1+t^2}\right),\qquad
 t=\frac{a_2}{1+a_1}.
\]
The denominator in the inverse is $2/(1+t^2)>0$.
Section~\ref{sec:transport} justifies using this normalization for
certificate existence without a Ramsey-invariance input.

Write the distinct finite parameters as $t_1,\ldots,t_n$ and let
$k=\trdeg_\Q F$. Choose a transcendence basis $u_1,\ldots,u_k$.
The field $F$ is finite separable over $\Q(u_1,\ldots,u_k)$.
Each partial derivation extends uniquely to $F$: for an algebraic
element $\beta$ with minimal polynomial $f$, differentiating
$f(\beta)=0$ prescribes
\[
 D(\beta)=-\frac{(Df)(\beta)}{f'(\beta)}.
\]
Here $Df$ differentiates the coefficients. Separability makes the
denominator nonzero, and the prescription preserves the defining
ideal of the simple extension. Applying this to finite simple
extensions proves existence and uniqueness. It follows that
$\Der_\Q(F,F)$ has basis $D_1,\ldots,D_k$, where
$D_j(u_l)=\delta_{jl}$. Its common constants are precisely the
elements of $F$ algebraic over $\Q$: a transcendental element can
itself be included in a transcendence basis and assigned derivative
one.

Let
\begin{equation}\label{eq:velocity-space}
 V=\bigl\{(D(t_i))_{i=1}^n:D\in\Der_\Q(F,F)\bigr\},\qquad
 W=\bigl\{(q(t_i))_{i=1}^n:\deg q\le2\bigr\}.
\end{equation}
Since the $t_i$ generate $F$, the velocity map is injective, and
$\dim_F V=k$. Its first three entries vanish. A quadratic vanishing
at those three distinct parameters is zero, so
\begin{equation}\label{eq:transversality}
 V\cap W=0,\qquad \dim W=3.
\end{equation}

Set
\begin{equation}\label{eq:barycentric}
 H(T)=\prod_{i=1}^n(T-t_i),\qquad \gamma_i=H'(t_i)^{-1}.
\end{equation}
Lagrange interpolation, followed by comparison of the coefficient of
$T^{n-1}$, gives
\begin{equation}\label{eq:moments}
 \sum_i\gamma_i t_i^r=0\quad(0\le r\le n-2),\qquad
 \sum_i\gamma_i t_i^{n-1}=1.
\end{equation}

\begin{proposition}[The graph condition]\label{prop:graph}
Suppose $n\ge7$ and the circle set has a tensor certificate. There is
an $F$-linear map $Q:V\to W$ such that
\begin{equation}\label{eq:graph-relation}
 B_l(v,Qw)+B_l(w,Qv)=4B_l(v,w)
 \quad(v,w\in V,\ 0\le l\le n-6),
\end{equation}
where $B_l(a,b)=\sum_i\gamma_i t_i^l a_ib_i$.
Consequently
\begin{equation}\label{eq:graph-U}
 U=\{v-\tfrac14Qv:v\in V\}
\end{equation}
has dimension $k$, is totally isotropic for every $B_l$, and satisfies
$V\subset U+W$.
\end{proposition}
\begin{proof}
Normalize the certificate by \eqref{eq:full-normalization}. In the
truncated ring
\[
 \cR=F[\delta_1,\ldots,\delta_k]/(\delta_1,\ldots,\delta_k)^3
\]
put $D_\delta=\sum_j\delta_jD_j$, acting trivially on the $\delta_j$.
The map
\[
 \varphi(x)=x+D_\delta x+\tfrac12D_\delta^2x
\]
is a $\Q$-algebra homomorphism. This follows by applying Leibniz
twice and discarding terms of degree at least three. In particular
it respects inverses, since a nonzero constant term is a unit.
Commutation of the individual derivations is not needed for this
calculation. Map the left factor of $\cB_F$ identically to $F$ and
the right factor by $\varphi$.

Using $h(1,t)$, the image matrix represents a polynomial
\[
 \cG(t,u)=-2(t-u)^2+\cG_1(t,u)+\cG_2(t,u),
\]
of degree at most two in each variable, with $\cG_j$ homogeneous
of degree $j$ in $\delta$. The first-order node equations give
$\cG_1(t_i,t_i)=0$. Each coefficient on the diagonal is a polynomial
of degree at most four with at least seven roots; hence
$\cG_1(t,t)=0$. Polynomial division gives
\[
 \cG_1(t,u)=(u-t)N(t,u),\qquad
 N(t,t)=\sum_j a_j(t)\delta_j,\quad \deg a_j\le2.
\]
The quotient $N$ has degree at most one in each variable: multiplication
by $u-t$ raises the degree in either variable by one. Put
$q(t)=\cG_2(t,t)$, a polynomial in $t$ of degree at most four.
With $v_i=\sum_jD_j(t_i)\delta_j$, the second-order node equations
are
\begin{equation}\label{eq:quadratic-node}
 q(t_i)=2v_i^2-v_i\sum_j a_j(t_i)\delta_j.
\end{equation}
The quadratic displacement in $\varphi(t_i)$ contributes no term of
order two: the leading kernel has a double zero on the diagonal,
and $\cG_1$ already has order one in $\delta$.

For a velocity $v=\sum_j c_j(D_j(t_i))_i$, define
$Qv=(\sum_jc_j a_j(t_i))_i$. This is well-defined because the
velocity columns are independent, and it takes values in $W$.
Multiplying \eqref{eq:quadratic-node} by $\gamma_i t_i^l$ and
summing, \eqref{eq:moments} gives
$B_l(v,Qv)=2B_l(v,v)$ for $0\le l\le n-6$.
Polarization proves \eqref{eq:graph-relation}.
Also \eqref{eq:moments} makes $W$ totally isotropic for every
such $B_l$, because $l+4\le n-2$. Expanding
$B_l(v-Qv/4,w-Qw/4)$ now gives zero. The graph map is injective
by \eqref{eq:transversality}, so $\dim U=k$; its definition gives
$V\subset U+W$.
\end{proof}

There is a useful equivalent form of this condition. The
$(n-5)\times n$ matrix $(\gamma_i t_i^l)_{0\le l\le n-6}$ has
rank $n-5$ by the Vandermonde determinant. Its kernel is exactly
the five-dimensional space of values of polynomials of degree at
most four, by \eqref{eq:moments}. Thus
\begin{equation}\label{eq:product-property}
 u,w\in U\quad\Longrightarrow\quad
 (u_iw_i)_{i=1}^n\text{ is the value vector of a polynomial of degree }\le4.
\end{equation}
This formulation will allow restriction to eight nodes.

\section{Common isotropic three-dimensional subspaces}\label{sec:subspaces}

The linear algebra in this section is over an arbitrary field of
characteristic zero. Nodes are distinct, $\gamma_i$ is defined by
\eqref{eq:barycentric}, and $W$ is the quadratic evaluation space.
A subspace is \emph{totally isotropic} for a symmetric bilinear form
if that form vanishes on every pair of its vectors.

\begin{lemma}[Seven nodes]\label{lem:seven}
For seven nodes, every three-dimensional subspace totally isotropic
for $B_0$ and $B_1$ is
\begin{equation}\label{eq:sign-space}
 U=E W,\qquad E=\diag(\eps_1,\ldots,\eps_7),\quad \eps_i\in\{1,-1\}.
\end{equation}
There are exactly $2^6$ such spaces.
\end{lemma}
\begin{proof}
Let $T=\diag(t_i)$. Since $B_1(a,b)=B_0(Ta,b)$, we have
$TU\subset U^{\perp_{B_0}}$. The latter has dimension four and
contains $U$. Thus
$K=\{u\in U:Tu\in U\}$ has dimension at least two. The linear map
$K\to U/K$, $u\mapsto Tu\bmod K$, has a nonzero kernel. Choose
$u\ne0$ in it; then $u,Tu,T^2u\in U$.

These vectors are independent. Otherwise $f(T)u=0$ for some nonzero
polynomial of degree at most two, so $u$ is supported on at most two
nodes. But $B_0(u,u)=B_1(u,u)=0$ is an invertible Vandermonde
system on those nodes for the quantities $\gamma_i u_i^2$. It
would force $u=0$.

The three vectors therefore form a basis of $U$. Applying both
forms to their pairs yields
\[
 \sum_i\gamma_i t_i^r u_i^2=0\qquad(0\le r\le5).
\]
The $6\times7$ Vandermonde matrix has a one-dimensional kernel,
spanned by $(\gamma_i)_i$ by \eqref{eq:moments}. Hence
$u_i^2=c$ for every $i$. Since $u\ne0$, $c\ne0$. Dividing by
$u_1$ gives $u_i/u_1=\eps_i\in\{1,-1\}$ and proves
\eqref{eq:sign-space}. Conversely every such space is isotropic
by \eqref{eq:moments}.

For the count, suppose a relative sign matrix preserves $W$.
Its sign vector is the evaluation of a polynomial $q$ of degree at
most two, and its product with $(t_i^2)_i$ is the evaluation of
another such polynomial $r$. The polynomial $T^2q(T)-r(T)$ has
degree at most four and seven roots, so is zero. Thus $q$ is
constant, and all relative signs coincide. Only a common sign
is redundant, giving $2^6$ spaces.
\end{proof}

\begin{lemma}[Eight nodes]\label{lem:eight}
For eight nodes, a three-dimensional subspace totally isotropic
for $B_0,B_1,B_2$ is $EW$ for a sign matrix $E$.
\end{lemma}
\begin{proof}
The map $T$ is injective on $U$. Indeed a vector in its kernel
would be supported at the unique possible zero node, and a
nonzero such vector has $B_0(u,u)\ne0$.
The subspace $U+TU$ is totally isotropic for $B_0$: the
$U,U$ pair uses $B_0$, the mixed pair uses $B_1$, and the
$TU,TU$ pair uses $B_2$. Nondegeneracy of $B_0$ in dimension
eight gives $\dim(U+TU)\le4$. Since $\dim TU=3$, the space
$K=\{u\in U:Tu\in U\}$ again has dimension at least two.
The preceding kernel argument supplies $u,Tu,T^2u\in U$;
their independence follows as in Lemma~\ref{lem:seven}.
The three forms on these vectors give
\[
 \sum_i\gamma_i t_i^r u_i^2=0\qquad(0\le r\le6).
\]
The $7\times8$ Vandermonde kernel is spanned by $(\gamma_i)_i$,
so all $u_i^2$ are equal and nonzero. Normalization gives $U=EW$.
The sign-uniqueness argument in Lemma~\ref{lem:seven} also
applies to eight nodes; there are $2^7$ such spaces.
\end{proof}

\begin{proposition}[Restriction and gluing]\label{prop:gluing}
Let $n\ge7$. Suppose $U\subset F^n$ has dimension three and
satisfies the product property \eqref{eq:product-property}.
Then $U=EW$ for a sign matrix $E$.
\end{proposition}
\begin{proof}
For $n=7$, the product property implies isotropy for $B_0,B_1$
by interpolation, so use Lemma~\ref{lem:seven}. Suppose $n\ge8$.
Choose eight coordinate positions such that restriction of $U$
has rank three: choose three independent coordinate rows of a
basis matrix and add five positions. The restricted product vectors
are quartic evaluations. Interpolation on these eight nodes makes
the restricted space isotropic for its own $B_0,B_1,B_2$.
Lemma~\ref{lem:eight} makes it a signed quadratic evaluation
space. In particular, any three of its coordinate rows are
independent.

For any additional node, combine it with any seven of the chosen
eight. The new restriction has rank three and is again a signed
quadratic space. On the common seven nodes the two spaces agree:
both are the image of $U$ there, and restriction has rank three.
The uniqueness of signs in Lemma~\ref{lem:seven} makes their
signs agree up to one common reversal. After this reversal, a
fixed vector of $U$ is represented on both sets by quadratic
polynomials agreeing at seven nodes; they are the same polynomial.
Thus the new sign extends the original signs consistently for
every vector in $U$. Repeating for each remaining node proves
$U=EW$ globally.
\end{proof}

\begin{lemma}[The two sign groups]\label{lem:sign-groups}
Suppose $n\ge7$, $\dim V=3$, $V\cap W=0$, and
$V\subset W+EW$. Then each sign occurs at least three times.
Moreover, restriction of every vector of $V$ to either sign group
is the evaluation of a quadratic polynomial on that group.
\end{lemma}
\begin{proof}
Let $m$ be the size of the smaller sign group $I$ and let $J$ be
the other group. If $m\le2$, then $|J|\ge5$.
For a vector in $W\cap EW$, write its values as $r(t_i)$ and
$\eps_iq(t_i)$ and reverse all signs so that $\eps_i=1$ on $J$.
On $J$ the quadratics $r,q$ agree, hence are identical. On $I$
the polynomial $q$ must vanish. Therefore
$\dim(W\cap EW)=3-m$ and
\[
 \dim((W+EW)/W)=m<3.
\]
This contradicts the injective image of $V$ in the quotient.
Both groups consequently have size at least three.
Finally, an element of $W+EW$ has the form
$(p(t_i)+\eps_iq(t_i))_i$. Its restrictions are the quadratics
$p+q$ and $p-q$, respectively.
\end{proof}

All these classifications are field arguments, rather than
enumerations of rational examples. As arithmetic checks,
\ancfile{structure.py} verifies the $64$ seven-node sign spaces
and their projection dimensions, and \ancfile{eight_factory.py}
checks the $128$ eight-node sign spaces and $\dim(U+TU)=4$
on distinct rational nodes. These checks supplement the proofs
above.

\section{Proof of the classification}\label{sec:classification}

\subsection{Necessity}
Suppose the circle set in Theorem~\ref{thm:main} is Ramsey.
Proposition~\ref{prop:necessity} gives a tensor certificate.
Transport and descent give a certificate for the finite rational-frame
normalization of Section~\ref{sec:second-order}. In its moduli field,
$\dim V=3$ and $V\cap W=0$. Proposition~\ref{prop:graph} provides
a three-dimensional space $U$ satisfying the product property,
with $V\subset U+W$. Proposition~\ref{prop:gluing} gives $U=EW$,
and Lemma~\ref{lem:sign-groups} partitions the nodes into two
groups of size at least three, on each of which every derivation
has quadratic parameter velocity.

For four distinct finite parameters in one group put
\[
 \rho=\frac{(t_a-t_c)(t_b-t_d)}{(t_a-t_d)(t_b-t_c)}.
\]
If $v_i=D(t_i)$, the quotient rule gives
\begin{equation}\label{eq:cross-derivative}
 \frac{D\rho}{\rho}=
 \frac{v_a-v_c}{t_a-t_c}+\frac{v_b-v_d}{t_b-t_d}
 -\frac{v_a-v_d}{t_a-t_d}-\frac{v_b-v_c}{t_b-t_c}.
\end{equation}
For quadratic velocities $v_i=a+bt_i+ct_i^2$, each difference
quotient is $b+c(t_i+t_j)$, so the four terms cancel. Thus every
$\Q$-derivation of $F$ kills $\rho$. The constants characterization
in Section~\ref{sec:second-order} makes $\rho$ algebraic over
$\Q$. All internal cross ratios in both groups are algebraic.
This proves the necessary implication using neither direction of
OpenAI's Theorem~1.1.

\subsection{The known sufficient mechanism}\label{sec:positive}
For completeness we record the circle certificate behind the
sufficient implication. This is a case of Strinz's two-group
theorem~\cite[Theorem C2]{strinz2026spherical}; the conclusion
that the certificate implies Ramsey is OpenAI's sufficient direction.

In the unit circle coordinate field $F$, choose
$A_I,A_J\in\mathrm{GL}_2(F)$ taking three points of each group to
$0,1,\infty$. The other normalized parameters are real algebraic
elements of $F$. Choose algebraic homogeneous representatives
$w_i$, so $A_Hz_i=\lambda_iw_i$ with $\lambda_i\in F^\times$
for $i\in H$, $H=I,J$.

There is an element $e\in\cB_F$ such that
\begin{equation}\label{eq:algebraic-projector}
 m(e)=1,\qquad e(L(\alpha)-R(\alpha))=0
\end{equation}
for every coordinate $\alpha$ of the finitely many $w_i$.
Here is an explicit construction. If $f_\alpha\in\Q[X]$ is the
minimal polynomial, put
\[
 g_\alpha(X)=\frac{f_\alpha(X)}{X-\alpha}\in F[X],\qquad
 e_\alpha=L(f_\alpha'(\alpha)^{-1})
                       g_\alpha^L(R(\alpha)),\qquad
 e=\prod_\alpha e_\alpha.
\]
The superscript $L$ embeds the coefficients in the left factor.
The identities
$(R(\alpha)-L(\alpha))g_\alpha^L(R(\alpha))=0$
and $g_\alpha(\alpha)=f_\alpha'(\alpha)\ne0$ prove
\eqref{eq:algebraic-projector}.

For $z=(s,t)^T$, $w=(u,v)^T$, define the bilinear graph forms
\[
 K_H(z,w)=\det\bigl(L(A_H)z,R(A_H)w\bigr)\qquad(H=I,J).
\]
At a node of group $H$, \eqref{eq:algebraic-projector} gives
$eK_H(L(z_i),R(z_i))=0$. Multiplication of coefficients gives
$m(K_H)=\det(A_H)(sv-tu)$. Hence
\begin{equation}\label{eq:two-group-certificate}
 G(z,w)=-2L((\det A_I\det A_J)^{-1})eK_I(z,w)K_J(z,w)
\end{equation}
vanishes at every tensor node, and $m(G)=-2(sv-tu)^2$.
The basis identity \eqref{eq:circle-kernel} yields a matrix with
$m(P)=\HH$. Dividing the node vectors by $s_i^2+t_i^2>0$
gives \eqref{eq:tensor-criterion}. All denominators are nonzero
elements of $F$ in pure tensor factors; no domain assumption is used.
At least three distinct circle points span the plane, so OpenAI's
Theorem~1.1, sufficient direction, applies. This proves
Theorem~\ref{thm:main} and hence Corollary~\ref{cor:seven}.

In particular, the construction applies to a group of four points
with algebraic cross ratio and any remaining triple, without a
dimension restriction. More generally it includes the algebraic-plus-
three conclusion of Strinz. Both are positive Ramsey consequences
of the sufficient direction of OpenAI's Theorem~1.1. The scripts
\ancfile{structure.py} and \ancfile{eight_factory.py} check the
diagonal and graph identities used by these constructions.

\section{The example in P\'alv\"olgyi's Appendix E}\label{sec:appendix-example}

\begin{proof}[Proof of Corollary~\ref{cor:palvolgyi}]
Write $w=-4-x-2y$. Algebraic independence makes the seven parameters
in \eqref{eq:appendix-family} distinct. The three rational anchors
$-2,0,2$ give a normalization defined over $\Q$ and with rational
inverse. Therefore the moduli field is $\Q(x,y,z)$, of
transcendence degree three.

We show that every four-point cross ratio is transcendental.
A cross ratio containing $z$ is a nonconstant fractional linear
function of $z$ over $\Q(x,y)$, so is transcendental over that field.
For a quadruple not containing $z$, there are three cases. Three
rational anchors and one of $x,y,w$ give a nonconstant fractional
linear function of that variable. Two rational anchors and two of
$x,y,w$ give a nonconstant function of two independent variables:
each pair among $x,y,w$ is algebraically independent, since its
affine change of variables from $(x,y)$ is invertible. Finally,
one rational anchor $a\in\{-2,0,2\}$ and all of $x,y,w$ have,
in one ordering, cross ratio
\begin{equation}\label{eq:appendix-cross}
 \rho=\frac{(a-y)(x-w)}{(a-w)(x-y)}
      =\frac{2(a-y)(x+y+2)}{(a+4+x+2y)(x-y)}.
\end{equation}
It is nonconstant: its denominator has $x^2$ coefficient one,
whereas its numerator has $x^2$ coefficient zero and $y^2$
coefficient $-2$. Thus the numerator cannot be a constant multiple
of the denominator. Other orderings give one of the six invertible
fractional linear transforms of $\rho$, and remain nonconstant.

The only elements of a pure rational function field over $\Q$
algebraic over $\Q$ are its rational constants. To check this
directly, a separable polynomial relation would make all partial
derivatives zero. For a reduced fraction $P/Q$, the equality
$\partial_j(P/Q)=0$ gives $P\mid\partial_jP$ when $P\ne0$;
degree in the $j$th variable forces $\partial_jP=0$, and then
$\partial_jQ=0$. Characteristic zero and repetition in all variables
make both constant. None of the preceding cross ratios is therefore
algebraic. Corollary~\ref{cor:seven}, necessary direction, proves
non-Ramsey-ness. That direction was established without either
direction of OpenAI's Theorem~1.1.
\end{proof}

\subsection{A direct finite obstruction}
There is also a short exact certificate that uses the second-order
node equations directly, independently of the common-subspace
classification. For this family the parameter Jacobian with respect
to $(x,y,z)$ is
\[
 J=\begin{pmatrix}
 0&0&0\\0&0&0\\0&0&0\\1&0&0\\0&1&0\\-1&-2&0\\0&0&1
 \end{pmatrix}.
\]
Writing $a_j(t)=\sum_{r=0}^2c_{jr}t^r$ in
\eqref{eq:quadratic-node}, the nine coefficients must satisfy
the following twelve linear equations:
\begin{equation}\label{eq:appendix-system}
 \sum_i\gamma_i t_i^l
   \bigl(J_{ij}a_k(t_i)+J_{ik}a_j(t_i)\bigr)
 =4\sum_i\gamma_i t_i^lJ_{ij}J_{ik},
 \quad l=0,1,\quad 1\le j\le k\le3.
\end{equation}
For $j=k$, this is twice the coefficient equation for
$\delta_j^2$; for $j<k$ it is the coefficient equation for
$\delta_j\delta_k$. Thus the factor four is uniform.

Order the coefficient columns as
$(c_{10},c_{11},c_{12},c_{20},c_{21},c_{22},c_{30},c_{31},c_{32})$,
and order rows first by $l$, then by
$(j,k)=(1,1),(1,2),(1,3),(2,2),(2,3),(3,3)$.
Let $M$ be the first ten rows of the augmented matrix. At the
expression test $(x,y,z)=(-3,-1,3)$, the seven nodes are distinct,
so every denominator $H'(t_i)$ is nonzero. The matrix $720M$ is
\[
\setlength{\arraycolsep}{3.5pt}
\begin{pmatrix}
-28&-36&-12&0&0&0&0&0&0&64\\
-15&-45&-15&-14&-18&-6&0&0&0&120\\
1&3&9&0&0&0&-14&-18&-6&0\\
0&0&0&-30&-90&-30&0&0&0&300\\
0&0&0&1&3&9&-15&-45&-15&0\\
0&0&0&0&0&0&2&6&18&4\\
-36&-12&-84&0&0&0&0&0&0&48\\
-45&-15&-45&-18&-6&-42&0&0&0&120\\
3&9&27&0&0&0&-18&-6&-42&0\\
0&0&0&-90&-30&-90&0&0&0&180
\end{pmatrix}.
\]
Fraction elimination and integer Bareiss elimination give
\begin{equation}\label{eq:appendix-determinant}
 \det(720M)=1223059046400000,\qquad
 \det M=\frac1{30611001600000}.
\end{equation}
Both calculations and a left-null-vector certificate are in
\ancfile{second_jet_certificate.py}. The nonzero test value proves
that this augmented minor is a nonzero rational function of
$x,y,z$. At every algebraically independent triple it stays
nonzero, since its nonzero numerator cannot vanish. Nine coefficient
columns cannot then span the right-hand side of the ten-row system.
This rules out a tensor certificate; Proposition~\ref{prop:necessity}
gives an avoiding colouring with $15^7$ colours in every dimension,
without using either direction of OpenAI's classification.

The rational test is a test of finitely many rational expressions,
not a specialization homomorphism from $\Q(x,y,z)$ to $\Q$.
In particular it makes no non-Ramsey assertion about the rational
seven-point configuration used in the determinant computation.

\subsection{Why the derivation-certificate question is different}
The discussion in \cite[Appendix~E, after Theorem~A.8]{palvolgyi2026}
concerns an identity
\[
 \sum_i\lambda_i\|D(a'_i)\|^2=C\ne0
\]
for one derivation $D$, fixed real scalar weights $\lambda_i$,
and every congruent copy in $\R^3$. The absence of such an identity
is a restriction on this particular non-Ramsey certificate method.
It does not imply tensor-certificate existence. The joint
second-order conditions \eqref{eq:appendix-system} involve mixed
coefficients from three independent parameter directions and are
inconsistent. Thus the determinant obstruction and the reported
failure of the single-derivation method have compatible conclusions.

\section{An elliptic boundary example}\label{sec:elliptic}

We prove Theorem~\ref{thm:elliptic} by rational identities. No
elliptic-curve group law is needed. Work first in the function field
\[
 F=\Q(x,y),\qquad y^2=x^3-2,
\]
where $x$ is transcendental. Define the rational map
\begin{equation}\label{eq:elliptic-map}
 \begin{aligned}
 x'&=\frac{\cA(x,y)}{(x-3)^2},
 &\cA(x,y)&=3x^2+9x-4-10y,\\
 y'&=\frac{\cB(x,y)}{(x-3)^3},
 &\cB(x,y)&=5x^3+45x^2-40-(27x+19)y.
 \end{aligned}
\end{equation}
Write $(x_i,y_i)$ for its iterates, starting with $(x_0,y_0)=(x,y)$.
The circle parameters will be $x_0,\ldots,x_6$.

\subsection{A finite rational test and real embeddings}
The polynomial identity
\begin{equation}\label{eq:curve-identity}
 \cB(x,y)^2-\cA(x,y)^3+2(x-3)^6
 =(y^2-x^3+2)(-171x^2-1674x+1000y+1561)
\end{equation}
shows that the map preserves the curve wherever defined. To check
that the required denominators are nonzero as functions, and the
iterates distinct and nonconstant, use the curve point
$P_0=(129/100,-383/1000)$. Eight successive first coordinates are
the following fractions, denoted $r_i$:
\begin{center}
\small
\begin{tabular}{@{}cl@{}}
\toprule
$i$&$r_i=x_i(P_0)$\\
\midrule
0&$\dfrac{129}{100}$\\[3pt]
1&$\dfrac{164323}{29241}$\\[3pt]
2&$\dfrac{2340922881}{58675600}$\\[3pt]
3&$\dfrac{307326105747363}{160280942564521}$\\[3pt]
4&$\dfrac{794845361623184880769}{513127310073606144900}$\\[3pt]
5&$\dfrac{49680317504529227786118937923}{3458519104702616679044719441}$\\[3pt]
6&$\dfrac{30037088724630450803382035538503505921}{3010683982898763071786842993779918400}$\\[3pt]
7&$\dfrac{182386897568483763089689099250940650872600619203}{127572396335305049740547038646345741798859364401}$\\[3pt]
\bottomrule
\end{tabular}
\end{center}
Exact cross multiplication gives
\begin{equation}\label{eq:orbit-order}
 r_0<r_7<r_4<r_3<r_1<r_6<r_5<r_2,
 \qquad r_i\ne3\quad(0\le i\le7).
\end{equation}
The fractions and comparisons are checked by
\ancfile{elliptic.py}. Each iterate is therefore defined as a
rational function, and the first seven $x_i$ are distinct. If
$P_1$ is the image of $P_0$, then $x_i(P_1)=r_{i+1}\ne r_i$
for $0\le i\le6$, so each $x_i$ is nonconstant.

The rational function $x^3-2$ is not a square in $\Q(x)$, by its
odd order at infinity. For every transcendental real
$X>\sqrt[3]{2}$ and either sign $Y=\pm\sqrt{X^3-2}$, substitution
defines an embedding of $F$ into $\R$. Indeed $\Q(x)\simeq\Q(X)$,
and a relation $Y\in\Q(X)$ would make $x^3-2$ a square after
clearing denominators. The quadratic field extension therefore
embeds injectively. Every nonzero denominator and difference of
functions just checked remains nonzero at $(X,Y)$.
The actual circle coordinate field is the embedded copy of $F$:
the coordinates recover $X_0=X$ and $X_1$, and
\[
 Y=\frac{3X^2+9X-4-(X-3)^2X_1}{10}.
\]

\subsection{The tensor ring in this example is a domain}
This property is needed only for the following construction. Put
$K=\Q(x,u)$. The element $x^3-2$ is not a square in $K$.
In $K(y)$, a proposed square root of $u^3-2$ has the form
$a+by$, $a,b\in K$. Comparing the $y$ coefficient in its square
gives $ab=0$. If $b=0$, it would make $u^3-2$ a square in $K$;
if $a=0$, it would make $(u^3-2)/(x^3-2)$ a square in $K$.
Both have odd order at $u=\infty$, so neither is possible.
Consequently $K(y,v)$ with $y^2=x^3-2$, $v^2=u^3-2$
has degree four over $K$.

The ring
\[
 A_0=\Q[x,u,y,v]/(y^2-x^3+2,v^2-u^3+2)
\]
is free over $\Q[x,u]$ with basis $1,y,v,yv$. These elements
are independent in the degree-four field just constructed, so
$A_0$ injects into it and is a domain. The tensor ring
$\cB_F=F\otimes_\Q F$ is the localization of $A_0$ at the
nonzero pure left and pure right elements. Thus it too is a domain.
This argument is specific to this function field; the preceding
classification and two-group certificate assumed no such property.

\subsection{A simultaneous invariant}
For the two copies of the function field write
\[
 Z=(x-u)^2,\qquad S=xu(x+u)-4+2yv,
 \qquad d=S/Z.
\]
The last quotient is taken only in the fraction field of $\cB_F$.
The same map \eqref{eq:elliptic-map} on both pairs $(x,y)$ and
$(u,v)$ preserves $d$. Here is a polynomial verification.
Write
\[
 D=x-3,\quad E=u-3,\quad a=\cA(x,y),\quad a'=\cA(u,v),
 \quad b=\cB(x,y),\quad b'=\cB(u,v),
\]
and put
\begin{equation}\label{eq:invariant-polynomial}
 \begin{split}
 M&=aa'(aE^2+a'D^2)-4D^4E^4+2bb'DE,\\
 \mathcal C&=M(x-u)^2-S(aE^2-a'D^2)^2.
 \end{split}
\end{equation}
Division by the two monic curve equations gives
\begin{equation}\label{eq:invariant-ideal}
 \mathcal C=Q_1(y^2-x^3+2)+Q_2(v^2-u^3+2).
\end{equation}
The quotients are displayed in Section~\ref{sec:certificates}, and
\ancfile{elliptic.py} performs the divisions with exact coefficients.
The transformed $S$ and $Z$ have common denominator $D^4E^4$,
with numerators $M$ and $(aE^2-a'D^2)^2$, respectively.
Thus \eqref{eq:invariant-ideal} proves invariance of $d$.

All the needed differences $L(x_i)-R(x_i)$ are nonzero. In fact
$x_i$ is nonconstant and hence not rational, so there is a
$\Q$-linear functional $\lambda:F\to\Q$ with
$\lambda(1)=0$ and $\lambda(x_i)=1$. Applying
$\lambda\otimes\mathrm{id}$ to a proposed equality
$x_i\otimes1=1\otimes x_i$ gives $1=0$.
The denominators in every $d_i$ are consequently nonzero.
The invariant gives $d_i=d$ along all seven simultaneous iterates.

Set
\[
 K(t,w)=tw(t+w)-4,\qquad C(t,w)=t^2w^2+8(t+w).
\]
The further polynomial identity
\begin{equation}\label{eq:biquadratic}
 K(t,w)^2-(t-w)^2C(t,w)=4(t^3-2)(w^3-2)
\end{equation}
implies, for $t_i=L(x_i)$ and $u_i=R(x_i)$,
\begin{equation}\label{eq:invariant-node}
 (t_i-u_i)^2d^2-2K(t_i,u_i)d+C(t_i,u_i)=0.
\end{equation}
Indeed $S_i=K(t_i,u_i)+2L(y_i)R(y_i)$ and
$d=S_i/(t_i-u_i)^2$; expanding \eqref{eq:biquadratic} proves
\eqref{eq:invariant-node}.

\subsection{A certificate with only pure-factor denominators}
In $\cB_F[t,w]$, define
\begin{equation}\label{eq:elliptic-certificate}
 G(t,w)=-2L((4y^2)^{-2})
 \bigl[(t-w)^2S^2-2K(t,w)ZS+C(t,w)Z^2\bigr].
\end{equation}
This has degree at most two in each variable. Its only inverse
is a pure left factor, since $4y^2\ne0$ in $F$.
Multiplying \eqref{eq:invariant-node} by $Z^2$ proves that
\eqref{eq:elliptic-certificate} vanishes at each tensor node,
first in the fraction field and therefore in the domain $\cB_F$.
On multiplying coefficients, $m(Z)=0$ and $m(S)=4y^2$, so
\[
 m(G(t,w))=-2(t-w)^2.
\]
Equation \eqref{eq:circle-kernel} gives a tensor matrix with
multiplied value $\HH$ and zero node evaluations. The seven points
are distinct, hence span the plane. OpenAI's Theorem~1.1,
sufficient direction, proves the positive Ramsey assertion of
Theorem~\ref{thm:elliptic}.

\subsection{All thirty-five cross ratios are transcendental}
For $0\le a<b<c<d\le6$, put
\[
 \rho_{abcd}=
 \frac{(x_a-x_c)(x_b-x_d)}{(x_a-x_d)(x_b-x_c)}\in F.
\]
At $P_0$ and $P_1$ its values are respectively
$\crat(r_a,r_b;r_c,r_d)$ and
$\crat(r_{a+1},r_{b+1};r_{c+1},r_{d+1})$, where here the
notation denotes the corresponding difference formula.
All differences in the denominators are nonzero by
\eqref{eq:orbit-order}. Computing first value minus second value,
the exact signs in lexicographic order of $(a,b,c,d)$ are
\begin{equation}\label{eq:cross-signs}
 \texttt{--+-++--++-----+--++++---++++---+++}.
\end{equation}
These are $35$ nonzero signs. The rational table, the difference
formula, and \eqref{eq:cross-signs} form a finite certificate;
\ancfile{elliptic.py} enumerates it with \ancfile{fractions.Fraction}
and saves every exact difference.

If $\rho_{abcd}$ were algebraic over $\Q$, its irreducible
polynomial would have a rational root upon evaluation at $P_0$.
This evaluation is legitimate for that finite expression: every
denominator used is nonzero at $P_0$. The irreducible polynomial
would therefore be linear and $\rho_{abcd}$ a rational constant,
contradicting its different value at $P_1$. Thus it is transcendental
in $F$, and injectivity of the real embedding preserves this fact.
All other orderings are invertible fractional linear transforms,
so all four-point cross ratios are transcendental.

Finally $M(A)\subset F$ and $\trdeg_\Q F=1$, while $M(A)$ contains
a transcendental cross ratio. Hence $\trdeg_\Q M(A)=1$.
This completes the proof of Theorem~\ref{thm:elliptic}, with the
positive Ramsey conclusion depending exactly on OpenAI's sufficient
direction as specified.

\section{Exact identity certificates}\label{sec:certificates}

The ancillary directory contains the five scripts cited in this note,
together with a validator that compares the certificates printed here
with their output. They use
Python~3, its standard library, and SymPy; all numerical tests are
exact rational or integer calculations. The README gives serial
commands with the resource bounds. In particular:
\begin{itemize}
\item \ancfile{identities.py} checks the homogeneous determinant
and circle-kernel identities and the barycentric moment formulas.
\item \ancfile{structure.py} checks the $64$ seven-node sign spaces,
their projection dimensions, and an algebraic-diagonal certificate.
\item \ancfile{eight_factory.py} checks the $128$ eight-node sign
spaces and the universal two-group graph identity.
\item \ancfile{second_jet_certificate.py} checks
\eqref{eq:appendix-determinant} by two determinant algorithms and
saves an exact left-null-vector certificate.
\item \ancfile{elliptic.py} checks the polynomial identities
\eqref{eq:curve-identity}, \eqref{eq:invariant-ideal}, and
\eqref{eq:biquadratic}; the rational orbit; and every cross-ratio
difference in \eqref{eq:cross-signs}. It also performs $21$ rational
tests of the simultaneous invariant.
\end{itemize}
Rational expression tests prove the specified nonvanishing and
nonconstancy claims. The general statements about all isotropic
spaces are proved in Section~\ref{sec:subspaces}, and the invariant
is proved by its polynomial identity, rather than by the $21$
sample tests.

\begingroup
\interdisplaylinepenalty=10000
For an explicit certificate of \eqref{eq:invariant-ideal}, the
two polynomial quotients are
\begin{align*}
Q_1=-20(u-3)^2\bigl(&5u^4x-15u^4+5u^3x^2-45u^3
 -54u^2vx+10u^2vy+112u^2v\\
&-45u^2x^2+135u^2x+54uvx^2-100uvx-60uvy
 +114uv\\
&-40ux+120u-20v^2x^2+120v^2x-180v^2
 -12vx^2\\
&-114vx+90vy+20x^2-180\bigr),
\end{align*}
and
\begin{align*}
Q_2=-20(x-3)^2\bigl(&-15u^2x^3-45u^2x^2+54u^2xy
 -12u^2y+60u^2\\
&+5ux^4+120ux^3-54ux^2y+135ux^2
 -100uxy-40ux\\
&-114uy-240u+10vx^2y-60vxy+90vy
 -15x^4\\
&-225x^3+112x^2y+114xy+120x+180\bigr).
\end{align*}
\endgroup
Substitution into \eqref{eq:invariant-ideal} and expansion gives
an identity in $\Q[x,u,y,v]$. The script obtains the same quotients
by successive polynomial division in $y$ and $v$ and verifies the
zero remainder.


\begingroup
\raggedright
\bibliographystyle{plain}
\bibliography{references}
\endgroup

\section*{Use of AI}
The proofs and the text of this note were produced with AI systems
(GPT-6.1 Sol, with editing by Claude Opus~5.5) under the author's
direction. They were checked by independent reviews with AI models
(GPT-6.1 Sol and Gemini~4 Argon) and, for the explicit identities and
finite arithmetic certificates, by the exact computations in the
ancillary files. No human expert has reviewed this note. The author
takes responsibility for the content. Every positive Ramsey conclusion
depends on the sufficient direction of OpenAI's classification,
Theorem~1.1 of result~172, at commit
\texttt{adc7f1241b42e322a6451854ab7e4b4c146bf78a}~\cite{openai2026}.
The negative conclusions are justified by the independent tensor-necessity
and colouring argument in Section~\ref{sec:tensors}, so do not require
either direction of that theorem.
\end{document}
